\documentclass[11pt]{article}
\usepackage[a4paper,margin=2.7cm]{geometry}
\usepackage{amsmath,amsthm,amssymb,stmaryrd}
\usepackage{mathrsfs}
\usepackage{tikz-cd} % commutative diagrams
\usepackage{enumitem} % compact lists
\usepackage{hyperref}

% ------------------------------------------------------------------
% Theorem environments
% ------------------------------------------------------------------
\newtheorem{theorem}{Theorem}[section]
\newtheorem{lemma}[theorem]{Lemma}
\newtheorem{corollary}[theorem]{Corollary}

\theoremstyle{definition}
\newtheorem{definition}[theorem]{Definition}
\newtheorem{remark}[theorem]{Remark}
\newtheorem{example}[theorem]{Example}

% ------------------------------------------------------------------
% TikZ styles
% ------------------------------------------------------------------
\tikzset{>=stealth}

% ------------------------------------------------------------------
% Convenience macros
% ------------------------------------------------------------------
\newcommand{\cat}[1]{\mathbf{#1}}
\newcommand{\Set}{\cat{Set}}
\newcommand{\Graph}{\cat{Graph}}
\newcommand{\Hom}{\operatorname{Hom}}
\newcommand{\id}{\mathrm{id}}

\title{Homework IV\\[0.4em]\large Typed graphs over a fixed type graph form a category}
\author{Lucian Dorin Crainic -- 1938430}
\date{}

\begin{document}
\maketitle
\vspace{-1.25em}

\section{Prerequisites \& notation}

We work with \emph{directed graphs}. A graph is a quadruple $G=(E,V,s,t)$ where
\[
  s,t: E \longrightarrow V
\]
are the source and target maps (so $s(e)$ is the start vertex of $e$, and $t(e)$ its end vertex).
A \emph{graph morphism} $f:G_1\to G_2$ between $G_1=(E_1,V_1,s_1,t_1)$ and $G_2=(E_2,V_2,s_2,t_2)$
is a pair of functions $f=(f_E,f_V)$ with $f_E:E_1\to E_2$ and $f_V:V_1\to V_2$ such that
\[
  s_2\circ f_E = f_V\circ s_1
  \qquad\text{and}\qquad
  t_2\circ f_E = f_V\circ t_1.
\]
Composition in $\Graph$ is componentwise:
$(g\circ f)_E := g_E\circ f_E$ and $(g\circ f)_V := g_V\circ f_V$, and identities are
$\id_G=(\id_{E},\id_{V})$.

\smallskip
Fix once and for all a \emph{type graph}
\[
  \mathsf{TG} = (E_{\mathsf{TG}},V_{\mathsf{TG}},s_{\mathsf{TG}},t_{\mathsf{TG}}).
\]
All typing data in what follows refer to this fixed graph.

\section{Typed graphs and their morphisms}

\begin{definition}[Typed graph]\label{def:typed-graph}
A \emph{typed graph over $\mathsf{TG}$} is a pair $(G,\mathsf{type})$ where
$\mathsf{type}:G\to \mathsf{TG}$ is a graph morphism.
Equivalently, $\mathsf{type}=(\mathsf{type}_E,\mathsf{type}_V)$ with
\[
  s_{\mathsf{TG}}\circ \mathsf{type}_E = \mathsf{type}_V\circ s,
  \qquad
  t_{\mathsf{TG}}\circ \mathsf{type}_E = \mathsf{type}_V\circ t.
\]
\end{definition}

\begin{definition}[Morphisms of typed graphs]\label{def:typed-morphism}
Given two typed graphs $(G_1,\mathsf{type}_1)$ and $(G_2,\mathsf{type}_2)$,
a \emph{typed-graph morphism}
\[
  f:(G_1,\mathsf{type}_1)\longrightarrow (G_2,\mathsf{type}_2)
\]
is a graph morphism $f:G_1\to G_2$ such that the \emph{typing triangle} commutes:
\[
  \mathsf{type}_2\circ f = \mathsf{type}_1,
\]
i.e.\ componentwise,
$\mathsf{type}_{2,E}\circ f_E = \mathsf{type}_{1,E}$ and
$\mathsf{type}_{2,V}\circ f_V = \mathsf{type}_{1,V}$.
\end{definition}

\section{Main result}

\begin{definition}[Category]\label{def:category}
A \emph{category} $\mathcal{C}$ consists of:
\begin{itemize}[nosep]
  \item a class of objects $\mathrm{Ob}(\mathcal{C})$;
  \item for each pair $A,B\in\mathrm{Ob}(\mathcal{C})$, a \emph{set} $\Hom_{\mathcal{C}}(A,B)$ of morphisms $A\to B$;
  \item for each triple $A,B,C$, a composition operation
  \[
    \circ:\ \Hom_{\mathcal{C}}(B,C)\times\Hom_{\mathcal{C}}(A,B)\to \Hom_{\mathcal{C}}(A,C);
  \]
  \item for each object $A$, an identity morphism $\id_A\in\Hom_{\mathcal{C}}(A,A)$;
\end{itemize}
satisfying the axioms:
\begin{itemize}[nosep]
  \item \textbf{Associativity:} for composable $f,g,h$ we have $(h\circ g)\circ f=h\circ(g\circ f)$;
  \item \textbf{Unit laws:} for $f:A\to B$, we have $\id_B\circ f=f=f\circ\id_A$.
\end{itemize}
\end{definition}

\begin{theorem}\label{thm:is-category}
The typed graphs over $\mathsf{TG}$ with typed-graph morphisms form a category,
denoted $\Graph_{/\mathsf{TG}}$.
\end{theorem}

\begin{proof}
We verify the five requested conditions one by one.

\begin{enumerate}[label=\textbf{(\arabic*)}, leftmargin=2.2em]

\item \emph{Class of objects.}
An object is a pair $(G,\mathsf{type})$ where $G$ is a graph and $\mathsf{type}:G\to\mathsf{TG}$ is a graph morphism.
This is well defined: for any fixed $G$, possible objects with underlying $G$ are in bijection with
$\Hom_{\Graph}(G,\mathsf{TG})$, and across all graphs $G$ we obtain a (possibly large) class of objects,
which is standard in category theory (large vs.\ small categories).

\item \emph{Set of morphisms between two objects.}
Fix two objects $(G_1,\mathsf{type}_1)$ and $(G_2,\mathsf{type}_2)$ with $G_i=(E_i,V_i,s_i,t_i)$.
A graph morphism $f:G_1\to G_2$ is a pair of functions $(f_E:E_1\to E_2,\ f_V:V_1\to V_2)$ satisfying
\[
  s_2\circ f_E = f_V\circ s_1,
  \qquad
  t_2\circ f_E = f_V\circ t_1.
\]
Such pairs form a set (a subset of the Cartesian product $E_2^{E_1}\times V_2^{V_1}$).
The additional \emph{typing} constraints
\[
  \mathsf{type}_{2,E}\circ f_E = \mathsf{type}_{1,E},
  \qquad
  \mathsf{type}_{2,V}\circ f_V = \mathsf{type}_{1,V}
\]
pick out a subset of those graph morphisms. Hence
\[
  \Hom_{\Graph_{/\mathsf{TG}}}\big((G_1,\mathsf{type}_1),(G_2,\mathsf{type}_2)\big)
  \;=\;
  \{\, f:G_1\to G_2 \text{ in }\Graph \mid \mathsf{type}_2\circ f=\mathsf{type}_1 \,\}
\]
is a \emph{set}.

\item \emph{Composition operation (well-defined and preserves typing).}
Let
\[
  f=(f_E,f_V):(G_1,\mathsf{type}_1)\to(G_2,\mathsf{type}_2),
  \qquad
  g=(g_E,g_V):(G_2,\mathsf{type}_2)\to(G_3,\mathsf{type}_3).
\]
Define $g\circ f:=(g_E\circ f_E,\ g_V\circ f_V)$.
\smallskip

\emph{(a) $g\circ f$ is a graph morphism.}
We check the source equation; the target one is analogous.
As functions $E_1\to V_3$,
\[
  \underbrace{(g_V\circ f_V)\circ s_1}_{\text{LHS}}
  \;=\;
  g_V\circ(\,f_V\circ s_1\,)
  \;=\;
  g_V\circ(\,s_2\circ f_E\,)
  \;=\;
  (\,g_V\circ s_2\,)\circ f_E
  \;=\;
  (\,s_3\circ g_E\,)\circ f_E
  \;=\;
  \underbrace{s_3\circ(\,g_E\circ f_E\,)}_{\text{RHS}},
\]
using (in order) associativity of function composition, the fact that $f$ preserves sources,
the fact that $g$ preserves sources, and associativity again.

\emph{(b) $g\circ f$ preserves typing.}
We must show $\mathsf{type}_3\circ(g\circ f)=\mathsf{type}_1$ as morphisms $G_1\to \mathsf{TG}$,
i.e.\ both on edges and on vertices:
\[
  \mathsf{type}_{3,E}\circ(g_E\circ f_E)
  \;=\;
  (\mathsf{type}_{3,E}\circ g_E)\circ f_E
  \;=\;
  \mathsf{type}_{2,E}\circ f_E
  \;=\;
  \mathsf{type}_{1,E},
\]
since $g$ is typed ($\mathsf{type}_3\circ g=\mathsf{type}_2$) and $f$ is typed
($\mathsf{type}_2\circ f=\mathsf{type}_1$). The same argument holds for vertices:
\[
  \mathsf{type}_{3,V}\circ(g_V\circ f_V)
  \;=\;
  (\mathsf{type}_{3,V}\circ g_V)\circ f_V
  \;=\;
  \mathsf{type}_{2,V}\circ f_V
  \;=\;
  \mathsf{type}_{1,V}.
\]
Hence $g\circ f$ is a well-typed morphism $(G_1,\mathsf{type}_1)\to(G_3,\mathsf{type}_3)$.

\item \emph{Identity morphisms (and unit laws).}
For any object $(G,\mathsf{type})$ with $G=(E,V,s,t)$,
the identity $\id_G=(\id_E,\id_V)$ in $\Graph$ is a graph morphism because
\[
  s\circ \id_E = s = \id_V\circ s,
  \qquad
  t\circ \id_E = t = \id_V\circ t.
\]
It is also \emph{typed} since
\[
  \mathsf{type}_E\circ \id_E = \mathsf{type}_E,
  \qquad
  \mathsf{type}_V\circ \id_V = \mathsf{type}_V,
\]
so $\mathsf{type}\circ \id_G=\mathsf{type}$.
Therefore $\id_G$ serves as the identity morphism $\id_{(G,\mathsf{type})}$.
Finally, for any $f:(G_1,\mathsf{type}_1)\to(G_2,\mathsf{type}_2)$,
\[
  \id_{(G_2,\mathsf{type}_2)}\circ f = f = f\circ \id_{(G_1,\mathsf{type}_1)}
\]
since composition is inherited componentwise from $\Graph$.

\item \emph{Associativity of composition.}
Given composable typed morphisms $f,g,h$, we have
\[
  (h\circ g)\circ f
  \;=\;
  \big((h_E\circ g_E)\circ f_E,\ (h_V\circ g_V)\circ f_V\big)
  \;=\;
  \big(h_E\circ(g_E\circ f_E),\ h_V\circ(g_V\circ f_V)\big)
  \;=\;
  h\circ(g\circ f),
\]
since function composition is associative on both components.
Thus composition in $\Graph_{/\mathsf{TG}}$ is associative.
\end{enumerate}

All five conditions hold, so $\Graph_{/\mathsf{TG}}$ is a category.
\end{proof}

\end{document}
